\section{Results}
\label{sec:results}

\todo{v10 results are not available. Everything below the pilot summary is
the planned analysis: hypotheses carried over from the pilot, not
outcomes. Replace each hypothesis paragraph with the measured result and
delete the ones the data does not address.}

\subsection{Pilot study}
\label{sec:pilot}
We first ran all methods on Wojtek over a wider level space (v8,
Table~\ref{tab:levelspace}, right column), with two additional ACCEL
variants and a DR variant with a random per-level difficulty (DR-mixed);
full numbers are in Appendix~\ref{app:pilot}. Four observations shaped the
benchmark:
\begin{itemize}
  \item \textbf{The starting distribution dominated.} Methods that start
    from easy levels (ACCEL, its variants and DR-easy) solved 7.7--8.2\% of
    uniform levels; the others (DR-mixed, \plr, PAIRED, DR) solved
    0.4--3.5\%.
  \item \textbf{Curation rescued a bad generator.} \plr\ and DR see the
    same uniform levels, but \plr\ solved 2.2\% vs.\ 0.4\% of uniform
    levels and 9.8\% vs.\ 1.0\% at $d=0.5$.
  \item \textbf{ACCEL matched DR-easy.} 7.7\% vs.\ 8.0\% uniform success,
    \cvar{0.25} 2.1\% vs.\ 2.2\%. ACCEL-SFL (learnability score plus grow
    mutation) gave the only gain over DR-easy, and only in the tail
    (\cvar{0.25} 2.8\%; mean 8.2\% vs.\ 8.0\%). ACCEL and
    \plr\ differ in both generator and mutation, so the pilot does not
    isolate the effect of mutation.
  \item \textbf{The space was too hard.} About 92\% of uniform levels were
    solved by no method, which compresses differences between the best
    methods and means DR trains mostly on levels it cannot learn from. This
    motivated the calibration of \S\ref{sec:levelspace}.
\end{itemize}

\subsection{Main comparison}
\label{sec:main}

\begin{table}[t]
  \centering\small
  \caption{Final evaluation on the v10 level space, both robots, mean over
  3 seeds. Uniform success with 95\% bootstrap CI, cross-fitted
  \cvar{0.1} and \cvar{0.25}, success at $d=0.5$, level reliability and
  held-out suite success.}
  \label{tab:main}
  \resulttable{main}{Main results table (\texttt{tables/main.tex}):
  rows = 6 methods $\times$ 2 robots; columns = uniform success [CI],
  \cvar{0.1}, \cvar{0.25}, $d=0.5$, $d=1.0$, level corr., suite success,
  solvable-subset success.}
\end{table}

\begin{figure}[t]
  \centering
  \placeholder{Learning curves: deterministic success on 512 fixed eval
  levels vs.\ environment steps, one panel per robot (and per
  $d\in\{0.5,1.0\}$), one line per method, mean over 3 seeds with min--max
  band. Source: \texttt{eval/d<d>/success\_rate} logged every 10
  iterations.}
  \caption{Learning curves.}
  \label{fig:curves}
\end{figure}

Hypotheses carried over from the pilot, to be tested on v10:
\begin{description}
  \item[H1 (generator).] Methods with an easy generator (DR-easy, ACCEL,
    ACCEL-SFL) outperform those with a uniform one (DR, \plr, PAIRED) on
    uniform success. With a calibrated space DR should no longer be near
    zero, so the gap is expected to shrink.
    \todo{result}
  \item[H2 (curation).] \plr\ improves over DR at the same generator.
    \todo{result; report the ratio with CIs on both robots}
  \item[H3 (mutation).] ACCEL improves over DR-easy, the only comparison
    that holds the generator fixed and adds curation and mutation.
    \todo{result}
  \item[H4 (score).] Relative to ACCEL, ACCEL-SFL gains more in the
    tail (\cvar{0.1} and \cvar{0.25}) than in mean success.
    \todo{result; confounded with the grow mutation
    (\S\ref{sec:methodlist})}
  \item[H5 (PAIRED).] PAIRED without distillation is the weakest UED
    method despite twice the simulation.
    \todo{result}
  \item[H6 (morphology).] The method ranking on Go1 matches Wojtek's.
    \todo{result; rank correlation across robots}
\end{description}

\subsection{Sim2sim transfer}
\label{sec:sim2sim}
\begin{table}[t]
  \centering\small
  \caption{Sim2sim transfer: success of MJX-trained policies in CPU
  MuJoCo under each physics variant, mean over 3 seeds.}
  \label{tab:sim2sim}
  \resulttable{sim2sim}{Sim2sim table (\texttt{tables/sim2sim.tex}):
  rows = methods $\times$ robots; columns = MJX reference,
  \texttt{mjx\_match}, \texttt{newton\_elliptic}, \texttt{half\_dt},
  \texttt{full\_collision}, \texttt{soft\_contact}.}
\end{table}
\begin{description}
  \item[H7 (solver sensitivity).] The closed-chain Wojtek loses more
    success than the open-chain Go1 when the solver or contact model
    changes, since its legs depend on the equality-constraint closure.
    \todo{result}
  \item[H8 (curriculum).] Curricula that train on harder terrain transfer
    better across physics variants.
    \todo{result; define ``harder'' via the composition logs of
    Figure~\ref{fig:composition}}
\end{description}

\subsection{What the curricula train on}
\begin{figure}[t]
  \centering
  \placeholder{Curriculum composition over training: per method, the mean
  physical value of each of the 8 level parameters among training levels
  (\texttt{level/<param>}) vs.\ environment steps, plus the genome distance
  from flat. One small multiple per parameter, one line per method; Wojtek
  top row, Go1 bottom row.}
  \caption{Curriculum composition.}
  \label{fig:composition}
\end{figure}
\todo{Describe which features each curriculum grows first (e.g.\ whether
ACCEL's buffer moves away from the easy generator at all, which it did not
on the earlier control-point space) once \texttt{level/<param>} is
logged.}

\subsection{Capability envelopes}
\begin{table}[t]
  \centering\small
  \caption{Capability envelopes: largest feature magnitude passed with
  50\% success, Wojtek units, mean over seeds.}
  \label{tab:envelopes}
  \resulttable{envelopes}{Envelope table (\texttt{tables/envelopes.tex}):
  rows = methods $\times$ robots; columns = rough (m), slope up/down
  ($^\circ$), stairs up/down (m), gaps (m), blocks (m).}
\end{table}

Per-feature envelopes are plotted in Figure~\ref{fig:envelopes}
(Appendix~\ref{app:extra}).
\todo{Which features separate methods; whether Go1 and Wojtek fail on the
same features (e.g.\ step-up for the small closed-chain robot).}
